\documentclass[a4paper,10pt]{report}\input{../0.Préambule/Préambule_cours_prof}\begin{document}

\pagestyle{empty}
\newpage
\noindent NOM, Prénom et classe : \dotfill

\section*{Nombres rationnels et pourcentages}
\addcontentsline{toc}{section}{Nombres rationnels et pourcentages}

\noindent \textit{Il sera tenu compte dans la notation de la propreté ainsi que de la justification apportée à chacune des réponses.}

\noindent \textit{Le barème est donné à titre indicatif. Il pourra être modifié ultérieurement.}

\noindent \textit{L'usage de la calculatrice est \textbf{interdite}.} \hfill \textit{Durée : 55 minutes}

\medskip
\paragraph{Exercice 1}: \hfill \textbf{6 points}

\begin{enumerate}
\item Écrire en chiffre le nombre douze millions trois cent vingt-deux mille cent vingt-deux.

\medskip\noindent\grandscarreaux{\linewidth}{0.8cm}
\item Dans ce nombre écrit en chiffres, quelle est :
\begin{enumerate}
\item la proportion de $1$ (on donnera le résultat sous forme décimale)

\medskip\noindent\grandscarreaux{\linewidth}{1.6cm}
\item la proportion de $2$ (on donnera le résultat sous forme fractionnaire)

\medskip\noindent\grandscarreaux{\linewidth}{1.6cm}
\item la proportion de $3$ (on donnera le résultat en pourcentage)
\end{enumerate}
\end{enumerate}

\noindent \begin{minipage}{13cm}
\medskip\hspace{17mm}\grandscarreaux{10.5cm}{1.6cm}
\begin{enumerate}
\setcounter{enumi}{2}
\item Dans le cercle ci-contre, colorier sa surface pour représenter :
\begin{enumerate}
\item en rouge, la proportion de $1$.
\item en vert, la proportion de $2$.
\item en bleu, la proportion de $3$.
\end{enumerate}
\end{enumerate}
\end{minipage}
\begin{minipage}{4cm}
\begin{center}
\begin{tikzpicture}[line cap=round,line join=round,>=triangle 45,x=1.0cm,y=1.0cm,scale=0.6]
\clip(-2.1,-2.1) rectangle (4.1,4.1);
\draw [line width=2.pt] (1.,1.) circle (3.cm);
\end{tikzpicture}
\end{center}
\end{minipage}


\medskip
\paragraph{Exercice 2}: \hfill \textbf{4 points}

François collectionne les voitures miniatures. un sixième de ses voitures sont en plastique, les autres sont en métal. Il possède $24$ voitures en plastique.
\begin{enumerate}
\item Combien François possède-t-il de voitures en tout ? Justifie en écrivant le calcul.

\medskip\noindent\grandscarreaux{\linewidth}{2.4cm}
\item Combien François possède-t-il de voitures en métal ? Justifie en écrivant le calcul.

\medskip\noindent\grandscarreaux{\linewidth}{2.4cm}
\end{enumerate}

\medskip
\paragraph{Exercice 3}: \hfill \textbf{5 points}

Quelle fraction suis-je ?
\begin{enumerate}
\item Mon dénominateur est le numérateur de $\dfrac{41}{17}$ et mon numérateur est le dénominateur de $\dfrac{53}{18}$

\medskip\noindent\grandscarreaux{\linewidth}{1.6cm}
\item Multiplié $7$, je permet d'obtenir $3$.

\medskip\noindent\grandscarreaux{\linewidth}{1.6cm}
\item J'ai le même numérateur que trois demis et la moitié du dénominateur de cinq douzièmes.

\medskip\noindent\grandscarreaux{\linewidth}{1.6cm}
\item En pourcentage, je donne $25\%$ et mon dénominateur vaut $1000$.

\medskip\noindent\grandscarreaux{\linewidth}{1.6cm}
\end{enumerate}

\medskip
\paragraph{Exercice 4}: \hfill \textbf{5 points}

Sur la demi droite graduée ci-dessous, placer les nombres : \quad $\dfrac{1}{2}$ \quad ; \quad $\dfrac{3}{2}$ \quad ; \quad $\dfrac{1}{4}$ \quad ; \quad $\dfrac{3}{4}$ \quad ; \quad $\dfrac{9}{8}$. 

\begin{center}
\begin{tikzpicture}[x=8cm]  
\draw[thick,->] (0,0) -- (2,0);
\foreach \x in {0,1}
\draw[shift={(\x,0)},color=black] (0pt,6pt) -- (0pt,-6pt) node [below] {\x};
\foreach \x in {0.125,0.25,0.375,0.5,0.625,0.75,0.875,1.125,1.25,1.375,1.5,1.625,1.75,1.875}	
\draw[shift={(\x,0)},color=black] (0pt,3pt) -- (0pt,-3pt) ;
\end{tikzpicture}
\end{center}
\vspace{5mm}

\medskip
\paragraph{Exercice bonus}: A faire uniquement lorsque vous avez fini le reste des exercices \hfill \textbf{1 points}

\medskip
Les instruments de mesure sont $3$ cruches de contenances respectives : $12$ litres, $7$ litres et $4$ litres.\medskip

\noindent La cruche de $12$ litres est pleines d'eau, les $2$ autres sont vides.\medskip

\noindent Indiquez des transvasements à l'aide des trois cruches uniquement, pour avoir exactement $6$ litres dans $2$ cruches.

\medskip\noindent\grandscarreaux{\linewidth}{5.6cm}
\end{document}